%!TEX root = ../main.tex


\makeatletter
\renewenvironment{abstract}{%
    \thispagestyle{empty}
    \@beginparpenalty\@lowpenalty
    \begin{center}%
        \bfseries \abstractname
        \@endparpenalty\@M
    \end{center}
}%

{\par\vfil\null}
\makeatother

\begin{abstract}
% Abstract in Deutsch
LocalLLM ist ein Agent, der Fragen zu privaten Rechnungen und Verträgen beantwortet. Er läuft vollständig auf einem Laptop mit 16\,GB Arbeitsspeicher und ohne Grafikkarte, damit keine persönlichen Daten den Rechner verlassen. Er liest PDF- und DOCX-Dateien mit ihrer Struktur ein: Tabellenzeilen mit Spaltenkopf, Kapitelpfad, Diagrammbeschreibung und geprüfte Kerndaten. Jede Antwort bekommt nummerierte Quellen. Ich vergleiche zwei Prototypen. P1 zeigt erst die Belege und lässt dann das Modell Werte mit Belegnummer liefern, die Code prüft. P2 ist ein ReAct-Agent, der selbst aus acht Werkzeugen wählt. Getestet wird nicht durch Ausprobieren, sondern mit begründeten Fragen, Sollwerten aus den Dokumenten und versiegelten Testsätzen. Ich nutze drei Kennzahlen: eine eigene Korrektheitsregel sowie Faithfulness und Context Recall aus RAGAS. Ein Reverse Engineering verfolgt eine Frage durch alle Bausteine und zeigt, wo Fehler entstehen. Auf 24 ungesehenen Fragen beantwortet P1 14 richtig, P2 7. Fragen nach Dokumenten, die es nicht gibt, lehnt P1 in drei von vier Fällen richtig ab, weil Code den Bestand prüft. Für P1 ist meist nicht die Suche der Engpass (Context Recall 0,80), sondern das kleine Sprachmodell beim Lesen der Belege; RAGAS-Faithfulness findet dabei unbelegte Aussagen, die eine reine Zahlprüfung übersieht.
\end{abstract}

\begin{otherlanguage}{english}
    \begin{abstract}
    % Abstract in Englisch
    LocalLLM is an agent that answers questions about private invoices and contracts. It runs entirely on a laptop with 16\,GB of RAM and no graphics card, so no personal data leaves the machine. It ingests PDF and DOCX files with their structure: table rows with column headers, section paths, chart descriptions and verified key facts. Every answer carries numbered sources. Two prototypes are compared. P1 shows the evidence first and then lets the model return values with evidence numbers, which code checks. P2 is a ReAct agent that chooses from eight tools itself. Testing does not rely on trying things out but on justified questions, expected values taken from the documents and sealed test sets. Three metrics are used: a rule-based correctness check plus faithfulness and context recall from RAGAS. A reverse-engineering walk-through follows one question through every component and shows where errors arise. On 24 unseen questions P1 answers 14 correctly, P2 7. P1 correctly refuses three of four questions about documents that do not exist, because code checks the inventory. For P1 the bottleneck is mostly not retrieval (context recall 0.80) but the small language model reading the evidence; RAGAS faithfulness catches unsupported claims that a pure number check misses.
\end{abstract}
\end{otherlanguage}
